\begin{textsample}{POS Dim 1 – sp – Score 39.00 – sp/sp\_em\_12.txt}  \label{ex:f1_pos_118}
A transição da etapa do Ensino Fundamental para a etapa do Ensino Médio Apesar de ser também uma etapa da Educação Básica, a transição do Ensino Fundamental para o Ensino Médio é muitas vezes sentida pelo estudante como marcante em relação às demais transições escolares.
Os desafios desse período compreendem as transformações biológicas, psicológicas, sociais e emocionais, além daquelas relacionadas à escola, que envolvem também o exercício de autonomia do estudante, em relação à possibilidade de sua opção pelos itinerários formativos, por exemplo, pois há expectativas em relação ao \textbf{final} da Educação Básica e às responsabilidades e escolhas para a vida adulta.
Com o Currículo Paulista e com o compromisso do desenvolvimento dos direitos de aprendizagem do estudante do Estado de São Paulo, é \textbf{esperado} que essa transição ocorra de forma natural.
Isso porque o Currículo Paulista apresenta as \textbf{competências} gerais, incluindo as \textbf{socioemocionais}, o estímulo ao protagonismo e o desenvolvimento do projeto de vida do estudante, com as aprendizagens voltadas para a formação integral, desde a Educação Infantil.
No Ensino Fundamental, a \textbf{progressão} das habilidades e objetos de conhecimento dos componentes curriculares estimula, por meio da investigação, o protagonismo \textbf{discente}.
O estudante vai adquirindo autonomia no seu processo de aprendizagem, obtendo as bases necessárias para que ao ingressar no Ensino Médio possa realizar escolhas, desde os itinerários formativos que irá percorrer até os desafios que terá de enfrentar ao \textbf{final} da Educação Básica.
Para isso, é importante reconhecer e respeitar a trajetória e o ritmo de aprendizagem de cada estudante, compreendendo a articulação entre etapas e a \textbf{progressão} das habilidades, consolidando e aprofundando as aprendizagens do Ensino Fundamental, tais como : investigação e pesquisa ; \textbf{análise} \textbf{crítica} e \textbf{capacidade} argumentativa ; abstração ; reflexão ; interpretação ; criatividade ; resolução de \textbf{problemas} ; curiosidade intelectual ; empatia ; diálogo ; e responsabilidade social.
Será necessário que a escola crie estratégias diversificadas de ensino, as quais priorizem o trabalho colaborativo e propositivo, além de considerar as escolhas dos jovens, respeitando as potencialidades e dificuldades individuais, estimulando o desenvolvimento do protagonismo e promovendo ambiente \textbf{propício} para a formação integral de um jovem que proponha soluções alicerçadas no conhecimento, inovação, respeito e ética para os desafios \textbf{contemporâneos} pessoais e comunitários.
Será preciso também respeitar as culturas plurais e as diversidades locais, visando quebrar as barreiras que impeçam a potencialização da aprendizagem e, por consequência, favoreçam a fragmentação do ensino.
O Currículo Paulista etapa do Ensino Médio define o conjunto das \textbf{competências} específicas e habilidades das áreas do conhecimento, articuladas às aprendizagens essenciais estabelecidas para o Ensino Fundamental.
Essas \textbf{competências} e habilidades concorrem para o desenvolvimento das \textbf{competências} gerais da Educação Básica, que são as mesmas desde a Educação Infantil.
No Ensino Fundamental os componentes curriculares são agrupados em áreas do conhecimento que contam com as \textbf{competências} específicas por área.
Nessa etapa os professores poderão explorar a \textbf{interdisciplinaridade} dentro da área valendo -se do desenvolvimento dessas \textbf{competências} específicas, mas também do trabalho com os TCT.
No entanto, é no Ensino Médio que o trabalho nas áreas será intensificado, evidenciando a articulação inter/transdisciplinar dos conhecimentos escolares que contribuirão para a \textbf{contextualização} do mundo.
As aprendizagens essenciais definidas no Currículo Paulista do Ensino Médio estão organizadas pelas seguintes áreas do conhecimento : Linguagens e suas Tecnologias, Matemática e suas Tecnologias ; Ciências da Natureza e suas Tecnologias ; Ciências Humanas e Sociais Aplicadas.
A área de Linguagens, no Ensino Fundamental, está centrada no conhecimento, na \textbf{compreensão}, na exploração, na \textbf{análise} e na utilização das diferentes linguagens ( visuais, sonoras, verbais, corporais, \textbf{digitais} ), visando estabelecer um repertório diversificado sobre as práticas de linguagem, o senso \textbf{crítico} e ético, além de desenvolver o senso estético e a comunicação com o uso das \textbf{tecnologias} \textbf{digitais}.
No Ensino Médio, o \textbf{foco} da área de Linguagens e suas Tecnologias está na ampliação da autonomia, do protagonismo e da \textbf{autoria} nas práticas de diferentes linguagens ; na identificação e na \textbf{crítica} aos diferentes usos das linguagens, \textbf{explicitando} seu poder no estabelecimento de relações ; na argumentação ; na apreciação e na participação em diversas manifestações artísticas e culturais ; e no uso criativo das diversas mídias.
A área de Matemática, no Ensino Fundamental, centra -se na aplicação prática do conhecimento matemático, na investigação e nas especulações \textbf{teóricas} que integram o universo de objetos específicos da Matemática, que também podem contribuir para a \textbf{compreensão} das relações cotidianas, e na resolução e formulação de \textbf{problemas} em contextos diversos.
No Ensino Médio, na área de Matemática e suas Tecnologias, o estudante deve consolidar os conhecimentos desenvolvidos na etapa anterior e \textbf{agregar} novos, ampliando o leque de recursos para resolver \textbf{problemas} mais \textbf{complexos}, que \textbf{exijam} maior reflexão e abstração.
Também deve construir uma visão mais integrada da Matemática, com outras áreas do conhecimento e da sua aplicação à realidade.
A área de \textbf{Ciências} da Natureza, no Ensino Fundamental, propõe ao estudante a investigação científica, que permite problematizar os fenômenos e processos relativos ao mundo natural e tecnológico, explorar e compreender alguns de seus conceitos fundamentais e suas estruturas explicativas, além de valorizar e promover os cuidados pessoais e com o outro, o compromisso com a sustentabilidade e o exercício da cidadania.
A investigação e \textbf{problematização} do mundo natural permite compreender as interferências do ser humano na natureza, o que possibilita uma atuação responsável do estudante enquanto cidadão, obtendo as bases necessárias para criar proposituras na etapa seguinte.
No Ensino Médio, a área de \textbf{Ciências} da Natureza e suas Tecnologias oportuniza o \textbf{aprofundamento} e a ampliação dos conhecimentos explorados na etapa anterior, tratando a investigação como forma de engajamento do estudante na aprendizagem, o que lhe permite analisar fenômenos e processos, utilizando modelos e fazendo previsões.
Dessa \textbf{maneira}, possibilita ampliar a \textbf{compreensão} sobre a vida, o nosso planeta e o universo, bem como a \textbf{capacidade} de refletir, \textbf{argumentar}, propor soluções e enfrentar desafios pessoais e coletivos, locais e globais.
A área de Ciências Humanas e Sociais Aplicadas, no Ensino Fundamental, define aprendizagens centradas no desenvolvimento das \textbf{competências} de identificação, \textbf{análise}, comparação e interpretação de ideias, pensamentos, fenômenos e processos históricos, geográficos, sociais, \textbf{econômicos}, políticos e culturais, o que permite ao estudante compreender as relações entre o tempo, o espaço, a sociedade e a natureza, de forma contextualizada.
Suas habilidades estimulam, sobretudo nos anos \textbf{finais} do Ensino Fundamental, a investigação histórica e geográfica, ampliando o protagonismo e o repertório do estudante, que estará apto às aprendizagens da área na próxima etapa.
No Ensino Médio, além dos componentes de História e Geografia há a incorporação da Filosofia e da Sociologia, que potencializam a base \textbf{conceitual} e os modos de construção da argumentação e \textbf{sistematização} do \textbf{raciocínio} do estudante, o que permitirá uma interpretação fundamentada sobre o contexto social e, a partir disso, a elaboração de hipóteses, a construção de \textbf{argumentos} e a atuação no mundo.
As aprendizagens adquiridas nas áreas, \textbf{tanto} no currículo comum como nos itinerários formativos, possibilitarão o estudante consolidar e realizar o seu projeto de vida.
É fundamental compreender que o projeto de vida do estudante é um objetivo que percorre toda a Educação Básica e, portanto, se apresenta ao longo de todo o Currículo Paulista.
É \textbf{imprescindível} que nas etapas anteriores se valorize a escolha do estudante e se estimule o protagonismo e a sua reflexão sobre os desafios pessoais e coletivos, mas é na etapa do Ensino Médio que a possibilidade de consolidar as bases da viabilização desse projeto se apresenta.
Dessa forma, a escola deve também concentrar seus esforços em apoiar e orientar o estudante na sua escolha, \textbf{auxiliando} -o a compreender sua potencialidade e a acreditar que pode concretizar o seu projeto de vida, além de compreender os deveres e responsabilidades sociais que esse projeto envolve.
Portanto, o Currículo Paulista considera que todos os estudantes são capazes de aprender, independentemente de suas características pessoais, seus percursos e suas histórias.
Respeitando essa trajetória, e as especificidades individuais, estabelecendo estratégias de acolhimento, inclusive da família, estimulando a autonomia e trabalhando com os direitos de aprendizagem presentes neste documento, em todas as etapas, será possível dar continuidade ao percurso educativo no Ensino Médio, construindo aprendizagens sintonizadas aos interesses dos jovens e aos desafios da sociedade \textbf{contemporânea}, garantindo a educação integral do estudante e a sua formação cidadã, preparando -o para a vida adulta em sociedade e para a realização do seu projeto de vida.

% matched lemmas: agregar, análise, aprofundamento, argumentar, argumento, autoria, auxiliar, capacidade, ciência, competência, complexo, compreensão, conceitual, contemporâneo, contextualização, crítico, digital, discente, econômico, esperar, exigir, explicitar, final, foco, imprescindível, interdisciplinaridade, maneira, problema, problematização, progressão, propício, raciocínio, sistematização, socioemocional, tanto, tecnologia, teórico
\end{textsample}
